\chapter{資産価格とファクター}
\label{ch:factor}

ヘッジファンドの戦略の多くは、「どの資産が、なぜ、平均して高いリターンを持つのか」という資産価格の理論の上に立っている。この章では、その理論と、それを測るための道具（ファクター・モデル）をまとめる。

\section{効率的市場の仮説}

Fama（1970）\cite{fama1970} は、市場の効率性を、価格に織り込まれている情報の範囲で三つに分けた。

\begin{itemize}
\item \textbf{ウィーク型}：過去の価格の情報は、すでに価格に入っている。過去の価格の動きから将来を予測しても儲からない。
\item \textbf{セミストロング型}：公開されたすべての情報（決算、ニュース）は、すでに価格に入っている。
\item \textbf{ストロング型}：公開されていない情報も、価格に入っている。
\end{itemize}

ただし、効率性の検定は、必ず「正しいリスクの調整の方法」の仮定と組みになる（結合仮説の問題）。ある戦略が平均して儲かっても、それが「市場の非効率」なのか「まだ測れていないリスクの対価」なのかは、データだけからは決まらない。

Grossman・Stiglitz（1980）\cite{grossman1980} は、市場が完全に効率的なら、情報を集める費用を回収できないので、誰も情報を集めなくなり、価格は効率的でなくなる、と指摘した。市場は「情報を集める人が費用を回収できる程度に」非効率なはずである。ヘッジファンドの存在は、この均衡の上にある。

\section{CAPM}

Sharpe（1964）\cite{sharpe1964}・Lintner（1965）\cite{lintner1965} の資本資産価格モデル（CAPM）は、すべての投資家が平均分散で最適化し、同じ見通しを持つとき、
\begin{equation}
\E[R_i] - R_f = \beta_i\,(\E[R_m] - R_f), \qquad \beta_i = \frac{\Cov(R_i, R_m)}{\Var(R_m)}
\end{equation}
が成り立つことを示した。個別の銘柄の平均リターンは、市場ポートフォリオとの共分散（ベータ）だけで決まり、銘柄固有の変動は分散で消せるので報われない。

CAPM の重要な帰結は、\textbf{時価総額加重の市場ポートフォリオが最適}ということである。LLM-BLM の Black--Litterman モデルが、時価総額加重の比率から「均衡のリターン」$\pi = \delta \Sigma w_{\text{mkt}}$ を逆算するのは、この考え方による（\chapref{ch:pc}）。

実証では、CAPM はうまく当てはまらない。ベータの高い株の平均リターンは CAPM の予測より低く、ベータの低い株は高い（Black・Jensen・Scholes 1972 以来の観察。第 II 部の低ボラティリティ）。

\section{マルチファクター・モデル}

\subsection{APT}

Ross（1976）\cite{ross1976} の裁定価格理論（APT）は、リターンがいくつかの共通のファクターで動くとき、
\begin{equation}
R_{i,t} = a_i + \sum_{k=1}^{K} b_{ik} f_{k,t} + \varepsilon_{i,t}
\end{equation}
裁定の機会がなければ、平均リターンはファクターへの感応度 $b_{ik}$ の線形関数になることを示した。ファクターが何かは、理論は教えない。

\subsection{Fama--French のモデル}

Fama・French（1993）\cite{ff1993} は、市場に加えて、\textbf{規模}（小型株から大型株を引いた SMB）と\textbf{割安}（簿価時価比率の高い株から低い株を引いた HML）の二つのファクターで、多くの株のポートフォリオの平均リターンが説明できることを示した（3 ファクター・モデル）。
\begin{equation}
R_{i,t} - R_{f,t} = \alpha_i + \beta_i\,\mathrm{MKT}_t + s_i\,\mathrm{SMB}_t + h_i\,\mathrm{HML}_t + \varepsilon_{i,t}
\end{equation}

その後、次のファクターが加えられた。

\begin{table}[H]
\centering\small
\caption{代表的なファクター・モデル}
\label{tab:factormodels}
\begin{tabular}{@{}P{0.33\linewidth}P{0.6\linewidth}@{}}
\toprule
モデル & ファクター \\
\midrule
CAPM & 市場 \\
Fama・French（1993） & 市場、規模（SMB）、割安（HML） \\
Carhart（1997）\cite{carhart1997} & ＋モメンタム（過去 1 年の勝ち組から負け組を引いた UMD） \\
Fama・French（2015）\cite{ff2015} & ＋収益性（RMW：収益性の高い株から低い株を引く）、投資（CMA：資産の伸びの小さい株から大きい株を引く） \\
Hou・Xue・Zhang（2015）\cite{hxz2015} & 市場、規模、投資、自己資本利益率（q ファクター・モデル） \\
\bottomrule
\end{tabular}
\end{table}

\subsection{ファクターの作り方}

ファクターは、ふつう\textbf{ポートフォリオの並べ替え}で作る。

\begin{enumerate}
\item 毎月（または毎年 6 月）、すべての銘柄を特性（簿価時価比率など）で並べる。
\item 上位 30\% と下位 30\%（または上位 10\% と下位 10\%）のポートフォリオを作る（時価総額加重か等加重）。
\item 「上位のリターン − 下位のリターン」を、その月のファクターのリターンとする。
\end{enumerate}

Fama・French は、規模と特性の 2 × 3 の並べ替えで、規模の影響を打ち消している。ファクターのリターンは、ロング・ショートのポートフォリオのリターンそのものなので、「ファクター投資」はこれを実際に売買することである。

\section{ファクターの検定}

\subsection{時系列の回帰}

ある戦略のリターンを、既知のファクターに回帰し、切片 $\alpha$ がゼロかを検定する。切片がゼロなら、戦略は既知のファクターの組み合わせで再現でき、新しい情報はない。多くのポートフォリオの切片が同時にゼロかを検定するのが、Gibbons・Ross・Shanken（1989）\cite{grs1989} の GRS 検定である。

\subsection{横断面の回帰}

Fama・MacBeth（1973）\cite{famamacbeth1973} の方法では、毎期 $t$ に、銘柄のリターンを特性（または推定したベータ）に横断面で回帰する。
\begin{equation}
R_{i,t+1} = \gamma_{0,t} + \gamma_{1,t}\,x_{i,t} + u_{i,t+1}
\end{equation}
係数 $\gamma_{1,t}$ の時系列の平均と標準誤差から、特性 $x$ がリターンを予測するかを検定する。毎期の係数は、特性 $x$ に比例して重みをつけたロング・ショートのポートフォリオのリターンと解釈できる。

\section{リスクか、ミスプライシングか}

ファクターの収益がなぜ生じるかについて、二つの説明がある。

\begin{table}[H]
\centering\small
\caption{ファクターの収益の二つの説明}
\label{tab:riskmisp}
\begin{tabular}{@{}P{0.15\linewidth}P{0.38\linewidth}P{0.38\linewidth}@{}}
\toprule
ファクター & リスクの説明 & 行動・ミスプライシングの説明 \\
\midrule
割安 & 割安な会社は業績が悪く、景気の悪いときにさらに損をする（苦境のリスク）。 & 投資家が過去の業績の悪さを将来に引き伸ばしすぎ、安く売りすぎる。 \\
モメンタム & 時間とともに変わるリスクへの感応度など、説明は難しい。 & 投資家が新しい情報に遅れて反応する（反応の不足）、または後で過剰に反応する。 \\
規模 & 小さい会社は流動性が低く、倒産しやすい。 & 注目されず、情報の反映が遅い。 \\
低ベータ & 借入れの制約のある投資家が、高ベータの株を好むため。 & 宝くじのような株を好む偏り。 \\
\bottomrule
\end{tabular}
\end{table}

どちらの説明でも、裁定には限界がある。Shleifer・Vishny（1997）\cite{shleifer1997} は、裁定をする運用者は他人の資金を使っていて、歪みが一時的に広がって損をすると、投資家が資金を引き上げるため、最も歪みの大きいときに裁定ができなくなることを示した（裁定の限界）。これは、LTCM や 2007 年の quant quake の説明にもなる（第 II 部）。

\section{ファクターの大きさ}

米国株式の長期のデータ（Kenneth French のデータ・ライブラリー、1920 年代後半から）で、代表的なファクターの月次の平均は、おおよそ次のとおりである（期間によって大きく変わる）。

\begin{itemize}
\item 市場の超過リターン：年 7〜8\% 前後、年率の標準偏差 18\% 前後。
\item 割安（HML）：年 3〜4\% 前後。ただし 2007 年ごろから 2020 年まで、長く損をした。
\item モメンタム（UMD）：年 7\% 前後と大きいが、まれに大きく崩れる（2009 年の春など）。
\item 規模（SMB）：年 2〜3\% 前後で、1980 年代以降は小さい。
\end{itemize}

どのファクターも、手数料・取引コストの前の値であり、何年も続けて損をする期間がある。ファクターの研究と、その崩れについては第 II 部（第 \ref{ch:equityfactor} 章）で扱う。

\section{ファクター・ズーと多重検定}

1990 年代以降、平均リターンを予測するとされる特性が数百も報告された。Cochrane（2011）\cite{cochrane2011} はこれを「ファクターの動物園」と呼び、どれが本当に独立な情報を持つのかを問題にした。Harvey・Liu・Zhu（2016）\cite{hlz2016} は、多くの研究者が同じデータで多くの特性を試していることを考えると、新しいファクターの $t$ 値は 2 ではなく 3 を超える必要があると論じた（\chapref{ch:backtest}）。

\section{ファクターとヘッジファンド}

ファクターは、ヘッジファンドの世界で三つの役割を持つ。

\begin{enumerate}
\item \textbf{戦略そのもの}：ファクターのロング・ショートは、クオンツ株式の中心である。
\item \textbf{リスクの物差し}：運用チームの収益のうち、ファクターで説明できる部分（安く手に入る部分）を除き、残りを「本当のアルファ」として評価する。マルチマネージャーは、運用チームにファクターのリスクを打ち消すよう求める。
\item \textbf{リスクモデル}：銘柄の共分散を、少数のファクターの共分散で近似する（\chapref{ch:risk}）。
\end{enumerate}

\begin{projbox}
LLM-BLM では、リッジ回帰の統制に、業種と時期の固定効果、規模（浮動株時価総額の対数）、発表の前のリターン（モメンタムに近い）を入れた。これは、「LLM の点数が、既知のファクターとは別の情報を持つか」を確かめるための、ファクター・モデルの考え方そのものである。
\end{projbox}
